\title{QLoRA: Efficient Finetuning of Quantized LLMs}

\begin{abstract}
We introduce \textsc{QLoRA}, a finetuning technique designed for optimal memory efficiency, allowing for the finetuning of large-scale models with up to 65B parameters on a single 48GB GPU while maintaining the performance characteristic of standard 16-bit finetuning. \textsc{QLoRA} achieves this by propagating gradients through a fixed, 4-bit quantized pretrained language model into Low Rank Adapters~(LoRA). The top-performing model series produced using this method, which we call \textbf{Guanaco}, achieves higher scores than all publicly released models to date on the Vicuna benchmark, attaining 99.3\% of ChatGPT's performance with only 24 hours of training on a single GPU. Several key techniques are proposed within \textsc{QLoRA} to enhance memory efficiency while retaining model effectiveness: (a) the introduction of 4-bit NormalFloat (NF4), a novel datatype proven to be information-theoretically ideal for normally distributed weights, (b) Double Quantization, which minimizes memory use by quantizing the quantization constants themselves, and (c) Paged Optimizers for mitigating episodic memory usage spikes. Employing \textsc{QLoRA}, we finetune over 1,000 models and comprehensively examine instruction-following and chatbot performance across 8 instruction datasets, a range of model architectures (LLaMA, T5), and model sizes that conventional finetuning cannot accommodate (including 33B and 65B parameter models). Our empirical findings demonstrate that using QLoRA for finetuning on modestly sized but high-quality datasets can surpass previous state-of-the-art results, even with smaller models. We further provide a nuanced evaluation of chatbot quality through both human assessments and GPT-4 based evaluations, highlighting the practicality and reliability of GPT-4 as an evaluation substitute. Additionally, we observe that existing chatbot benchmarks often fail to provide trustworthy appraisals of chatbot capabilities. Through targeted analysis of failure cases ("lemon-picked" examples), we identify situations where \textbf{Guanaco} does not match ChatGPT's performance. All models, codebase, and the necessary CUDA kernels for 4-bit training are openly released.\footnote{\url{https://github.com/artidoro/qlora} and \url{https://github.com/TimDettmers/bitsandbytes}}
\end{abstract}

\section{Introduction}

Adapting large language models (LLMs) through finetuning is a powerful approach for enhancing their capabilities \citep{min2021metaicl, wei2021finetuned, ouyang2022training, wang2022super, wang2022self, liu2022few}, as well as for incorporating or mitigating particular behaviors \citep{ouyang2022training,askell2021general,bai2022training}. Yet, finetuning extremely large models incurs significant computational costs; for instance, standard 16-bit finetuning of a LLaMA model with 65B parameters~\citep{touvron2023llama} necessitates over 780 GB of GPU memory. While advances in quantization can effectively lower the memory requirements for LLMs during inference \citep{dettmers2022llm, dettmers2022case, frantar2022gptq, xiao2022smoothquant}, these memory-saving strategies do not extend to the training phase, where they often fail \citep{wortsman2023stable}.

In this work, we present the first successful approach for finetuning a 4-bit quantized model without any loss in performance. Our approach, termed \textsc{QLoRA}, introduces a novel, high-precision quantization method that compresses a pretrained model to 4-bit. Subsequently, a lightweight set of trainable Low-rank Adapter parameters \citep{hu2021lora} is incorporated, 

which are optimized by propagating gradients through the quantized weights.

With \textsc{QLoRA}, the memory consumption necessary for finetuning a 65B parameter model drops from more than 780GB to under 48GB of GPU memory, all without sacrificing training speed or model accuracy when compared to the 16-bit fully finetuned counterpart. This improvement greatly democratizes access to LLM finetuning, making it feasible to finetune the largest models currently available using just a single GPU. Leveraging \textsc{QLoRA}, we introduce the \textbf{Guanaco} model suite: our second-best model attains 97.8\% of ChatGPT’s performance on the Vicuna~\citep{vicuna2023} benchmark, and can be trained in under 12 hours on a typical consumer GPU. Training our largest model for 24 hours on a single professional GPU achieves 99.3\% of ChatGPT’s performance on Vicuna, essentially eliminating the performance gap. Moreover, in deployment, our smallest \textbf{Guanaco} model (7B parameters) requires only 5 GB of memory while surpassing the performance of a 26 GB Alpaca model by over 20 percentage points on the Vicuna benchmark (see Table~\ref{tbl:vicuna_eval}).

\textsc{QLoRA} brings forth several key innovations aimed at reducing memory requirements while preserving model performance: (1) \textbf{4-bit NormalFloat}, a quantization datatype optimized from an information theoretic standpoint for data with normal distributions, which demonstrates superior empirical results over both 4-bit Integers and 4-bit Floats. (2) \textbf{Double Quantization}, which involves an additional quantization step for the quantization coefficients themselves, resulting in an average savings of roughly 0.37 bits per parameter (equivalent to about 3 GB for a model with 65B parameters). (3) \textbf{Paged Optimizers}, leveraging NVIDIA unified memory to mitigate memory spikes commonly caused by gradient checkpointing during mini-batch processing, especially with long input sequences. Collectively, these techniques are integrated into an enhanced LoRA method that places adapters at every layer of the network, thus circumventing most accuracy compromises noted in earlier studies.

Thanks to \textsc{QLoRA}'s efficiency, we are able to extensively explore instruction finetuning and chatbot capabilities at model scales previously inaccessible with standard finetuning, primarily due to prohibitive memory demands. Specifically, we train over 1,000 models spanning various instruction tuning datasets, model types, and parameter counts ranging from 80M to 65B. Our results confirm that \textsc{QLoRA} can match 16-bit finetuning performance (\S\ref{sec:comp_to_std}), and also facilitate the development of a state-of-the-art chatbot, \textbf{Guanaco} (\S\ref{sec:pushing}), along with a detailed examination of patterns across these trained models. Notably, our analysis highlights that data quality has a much greater impact than sheer dataset size. For example, a dataset of 9,000 samples (OASST1) outperformed a much larger 450,000 sample dataset (FLAN v2, subsampled) in chatbot benchmarks, even when both datasets are designed to enhance instruction following. Additionally, our findings reveal that high performance on the Massive Multitask Language Understanding (MMLU) benchmark does not necessarily predict success on the Vicuna chatbot benchmark, and vice versa, underscoring that the alignment between the dataset and the target task is more critical than the dataset's magnitude.

In addition, we conduct a thorough assessment of chatbot efficacy utilizing both human annotators and GPT-4 as evaluators. Our approach involves a tournament-based evaluation framework, in which models are pitted against one another in head-to-head matches to generate the most appropriate response for a specific prompt. Each match's outcome is decided by either GPT-4 or a human rater. We then summarize these head-to-head results using Elo scores~\citep{elo1967proposed, elo1978rating}, which serve to rank the chatbots according to their performance. Notably, the rankings produced by GPT-4 and human reviewers show high correspondence, although there are also instances where the two evaluators diverge significantly in their judgments. This observation underscores the utility of model-based evaluation as a lower-cost substitute for manual annotation, while also emphasizing its inherent limitations and potential unreliability.

Complementing our quantitative benchmarking, we provide a qualitative assessment of the \textbf{Guanaco} models. This qualitative exploration brings to light cases of both success and failure that quantitative results did not capture.

To support ongoing research, we make available all model responses along with their corresponding human and GPT-4 evaluations. Our codebase, CUDA kernels, and associated methods have been open-sourced and are integrated with the Hugging Face transformers ecosystem \citep{wolf2019huggingface}, ensuring broad accessibility. Additionally, we are releasing adapter modules for model sizes of 7/13/33/65B, each finetuned on 8 distinct instruction-following datasets, culminating in a set of 32 open-source, specialized models.

\section{Background}
\paragraph{Block-wise k-bit Quantization} Quantization refers to transforming a data representation with a higher informational resolution into one with lower precision, commonly by reducing the number of bits per element (e.g., converting 32-bit floating-point values to 8-bit integers). To effectively leverage the representational capacity of the target, low-bit format, the source data is typically normalized—by dividing by the largest absolute value within a tensor—before conversion. For example, converting a FP32 tensor to an Int8 tensor with values within $[-127, 127]$ proceeds as follows:

\begin{equation}
    \mathbf{X}^{ \text{Int8}} = \text{round}\left( \frac{127}{{\text{absmax}}(\mathbf{X}^{{\text{FP32}}})}\mathbf{X}^{{\text{FP32}}} \right) = \text{round}(c^{\text{FP32}}\cdot \mathbf{X}^{{\text{FP32}}}),
\end{equation}

where the parameter $c$ is known as the {\em quantization constant} or {\em quantization scale}. The reverse transformation—dequantization—recovers the original precision as follows:

\begin{equation}
    \text{dequant}(c^{\text{FP32}}, \mathbf{X}^{\text{Int8}}) = \frac{\mathbf{X}^{\text{Int8}}}{c^{\text{FP32}}} = \mathbf{X}^{\text{FP32}}
\end{equation}

This strategy faces a significant drawback: when the input tensor contains a value with an exceptionally large magnitude—an outlier—the quantization process becomes inefficient. Specifically, certain quantization bins (specific bit patterns) remain scarcely populated or entirely unused. To alleviate the effects of outliers, a widely adopted solution is to partition the input tensor into several blocks, each quantized separately and assigned its own scaling factor $c$. This method can be expressed as follows: The input tensor $\mathbf{X}\in \mathbb{R}^{b\times h}$ is divided into $n$ consecutive blocks, each of size $B$, by first flattening the tensor and then segmenting it linearly into $n = ({b\times h})/{B}$ distinct blocks. Each block undergoes independent quantization according to Equation~1, resulting in a quantized tensor and a unique quantization constant $c_i$ for every block.

\paragraph{Low-rank Adapters} The Low-rank Adapter (LoRA) technique for finetuning \citep{hu2021lora} decreases memory consumption by introducing a limited set of trainable adapter parameters, while keeping the primary model weights fixed. During stochastic gradient descent, gradients flow through the immutable pretrained model weights to the adapter, which is the only component being optimized to minimize the loss. In LoRA, a standard linear transformation is enriched with an additional, low-rank decomposition. If the baseline projection is ${\mathbf{X} \mathbf{W} = \mathbf{Y}}$ where $\mathbf{X}\in  \mathbb{R}^{b\times h}$ and $\mathbf{W}\in  \mathbb{R}^{h\times o}$, LoRA introduces:

\begin{equation}
    \mathbf{Y} = \mathbf{X}\mathbf{W} + s\mathbf{X}\mathbf{L}_1\mathbf{L}_2,
\end{equation}

Here, $\mathbf{L}_1\in  \mathbb{R}^{h\times r}$ and $\mathbf{L}_2\in  \mathbb{R}^{r\times o}$, and $s$ denotes a scalar scaling parameter.

\paragraph{Memory Requirement of Parameter-Efficient Finetuning} The memory consumption of LoRA during training warrants careful consideration, particularly regarding both the quantity and scale of adapters involved. Given that LoRA’s memory usage remains extremely low, it is feasible to deploy additional adapters for performance gains without a significant uptick in overall memory requirements. Although LoRA was introduced as a Parameter Efficient Finetuning (PEFT) technique, it is not the LoRA parameters themselves but the activation gradients that dominate memory consumption during LLM finetuning. For example, finetuning a 7B LLaMA model on FLAN v2 with a batch size of 1 and LoRA parameters corresponding to the widely adopted 0.2\% of the original model’s weights~\citep{hu2021lora,liu2022few} results in input gradients occupying 567 MB of memory, whereas the LoRA parameters themselves require merely 26 MB. Employing gradient checkpointing~\citep{chen2016training} reduces input gradient usage further, averaging 18 MB per sequence, yet this still exceeds the memory demanded by all LoRA weights combined. In stark contrast, the 4-bit quantized base model uses 5,048 MB of memory. These observations underscore not only the utility of gradient checkpointing but also that further minimizing LoRA parameters yields limited memory savings. Therefore, increasing the number of adapters does not appreciably impact the overall memory usage during training (refer to Appendix~\ref{app:memory} for a comprehensive breakdown). As will be elaborated upon, this is a key consideration for achieving performance on par with full 16-bit precision.

\section{\textsc{QLoRA} Finetuning}

\textsc{QLoRA} enables accurate 4-bit finetuning by leveraging two newly introduced methods—4-bit NormalFloat (NF4) quantization and Double Quantization. Furthermore, we present Paged Optimizers to mitigate sudden memory usage increases during gradient checkpointing, which have typically resulted in out-of-memory issues and hindered large model finetuning on a single device.

\textsc{QLoRA} utilizes a single low-precision data format for storage, generally at 4 bits, while employing a separate computation data type, usually BFloat16. In practice, this means that each time a \textsc{QLoRA} weight tensor is accessed, it is converted from the quantized form to BFloat16 prior to conducting 16-bit matrix multiplication.

Below, we outline the main elements of \textsc{QLoRA} and then present a rigorous definition of the \textsc{QLoRA} method.

\paragraph{4-bit NormalFloat Quantization} The NormalFloat (NF) format expands upon Quantile Quantization \citep{dettmers2022optimizers}, an information-theoretically optimal technique where every quantization interval (or "bin") contains the same number of mapped input tensor elements. This quantization method involves estimating tensor quantiles through the empirical cumulative distribution function.

A primary drawback of quantile quantization is its computational cost, as accurate quantile computation is resource intensive. Consequently, rapid quantile estimation schemes, such as SRAM quantiles~\citep{dettmers2022optimizers}, are adopted in practice. However, these fast algorithms introduce quantization inaccuracies, particularly affecting outliers, which are frequently among the most significant values.

These computational overheads and approximation-induced errors are minimized when the input tensors originate from a distribution that is invariant up to a quantization constant. In these scenarios, tensors exhibit identical quantiles, making precise quantile computation practical.

Given that pretrained neural network weights generally follow a normal distribution centered at zero with a standard deviation $\sigma$ (refer to Appendix~\ref{app:normality}), it is possible to map all weights to a common target distribution by adjusting $\sigma$ such that this distribution aligns precisely with the permissible bounds of our chosen data type. In this work, we define the data type range as $[-1, 1]$. Consequently, both the quantiles associated with the data type and those derived from the network weights must be scaled to occupy this specified interval.

To achieve the information-theoretically optimal quantization for zero-mean normal distributions with varying standard deviations $\sigma$ over $[-1, 1]$, we employ the following approach: (1) determine the $2^k+1$ quantile points from a standard $N(0,1)$ distribution, yielding a $k$-bit quantile-based quantization data type tailored for normal distributions; (2) rescale the resulting quantile levels to reside within $[-1, 1]$; (3) normalize an input weight tensor to $[-1, 1]$ by dividing by its absolute maximum. 

With the domains of the data and the quantization type aligned, the quantization procedure becomes straightforward. Notably, step (3) effectively rescales the weight tensor’s standard deviation so that it matches that of the constructed $k$-bit data type. Formally, the set of $2^k$ data type quantization points $q_i$ is computed as follows:

\begin{equation}
q_i  = \frac{1}{2}\left( Q_X\left(\frac{i}{2^k+1}\right) + Q_X\left(\frac{i+1}{2^k+1}\right)\right),
\end{equation}

where $Q_X(\cdot)$ denotes the quantile function associated with the standard normal distribution $N(0,1)$. One challenge inherent to symmetric $k$-bit quantization is the absence of an explicit encoding for the value zero. This omission poses difficulties, particularly when it is essential to represent zero-valued elements—such as padding—without error. To guarantee the presence of a discrete zero point (i.e., $0$) while fully utilizing all $2^k$ available quantization levels, we adopt an asymmetric quantization scheme. In this approach, we estimate the quantiles $q_i$ for two separate ranges: $2^{k-1}$ quantiles corresponding to negative values, and $2^{k-1}+1$ quantiles for positive values. These two sets of quantiles are then combined, after which one of the two resulting zero values—present in both ranges—is eliminated. The data type defined by this procedure yields quantization bins containing, in expectation, equal numbers of values. We refer to this construction as \emph{k-bit NormalFloat} (NFk), emphasizing its optimality from an information-theoretic perspective for distributions centered at zero and following a normal distribution. Details concerning the exact values constituting this data type appear in Appendix~\ref{app:nf4}.

\paragraph{Double Quantization{}} 
We propose a method termed \emph{Double Quantization} (DQ), wherein the quantization parameters themselves are subjected to further quantization in order to achieve greater memory efficiency. While a small block size is necessary to attain high-precision 4-bit quantization~\citep{dettmers2022case}, this also introduces significant memory overhead. For instance, with 32-bit quantization constants and a block size of 64 for $\mathbf{W}$, each parameter incurs an additional $32/64 = 0.5$ bits on average due to these constants. Double Quantization addresses this issue by reducing the memory overhead of quantization parameters.

More concretely, Double Quantization treats the initial set of quantization constants, $c_2^{\text{FP32}}$, as inputs to a secondary quantization phase. In this phase, the constants are themselves quantized, yielding $c_2^{\text{FP8}}$, along with a new set of higher-level quantization constants $c_1^{\text{FP32}}$. For this secondary quantization, we employ 8-bit floating-point representations and increase the block size to 256, as empirical results indicate that such 8-bit quantization incurs no loss in performance~\citep{dettmers2022case}. 
Given that the $c_2^{\text{FP32}}$ are strictly positive, we first subtract their mean value to center the data, thereby making symmetric quantization feasible. This method leads to a notable reduction in memory cost: using a block size of 64, the memory overhead per parameter decreases from $32/64=0.5$ bits to $8/64 + 32/(64\cdot256) = 0.127$ bits. This translates into a per-parameter savings of 0.373 bits.

\paragraph{Paged Optimizers} leverage NVIDIA’s unified memory~\footnote{\tiny \url{https://docs.nvidia.com/cuda/cuda-c-programming-guide}} capability, which automatically handles memory paging between the CPU and GPU at the granularity of individual pages. This functionality ensures that GPU computations remain correct even when the GPU memory is occasionally exhausted. Unified memory mimics conventional operating system paging between system RAM and secondary storage. In our implementation, optimizer states are allocated as paged memory so that, during GPU memory pressure, they are transparently evicted to CPU RAM and later brought back to the GPU as needed for the optimizer update.

\paragraph{\textsc{QLoRA}.} Building upon the mechanisms outlined above, we formalize \textsc{QLoRA} in the context of a quantized base model’s single linear layer enhanced by a single LoRA adapter as:

\begin{equation}
    \mathbf{Y}^{\text{BF16}} =  \mathbf{X}^{\text{BF16}} \text{doubleDequant}(c_1^{\text{FP32}}, c_2^{\text{k-bit}}, \mathbf{W}^{\text{NF4}}) + \mathbf{X}^{\text{BF16}}\mathbf{L}^{\text{BF16}}_1\mathbf{L}^{\text{BF16}}_2,
\end{equation}

The operator doubleDequant$(\cdot)$ is described by:

\begin{equation}
    \text{doubleDequant}(c_1^{\text{FP32}}, c_2^{\text{k-bit}}, \mathbf{W}^{\text{k-bit}}) = \text{dequant}(\text{dequant}(c_1^{\text{FP32}}, c_2^{\text{k-bit}}), \mathbf{W}^{\text{4bit}}) = \mathbf{W}^{\text{BF16}},
\end{equation}

Here, $\mathbf{W}$ utilizes the NF4 quantization scheme, while $c_2$ is encoded in FP8.

A block size of 64 is employed for $\mathbf{W}$ to improve quantization fidelity, whereas $c_2$ adopts a block size of 256 to achieve better memory efficiency.

When performing parameter updates, only the gradients of the loss with respect to the adapter weights, $\frac{\partial E}{\partial \mathbf{L}_i}$, are necessary; gradients for the 4-bit weights, $\frac{\partial E}{\partial \mathbf{W}}$, are not required. Nevertheless, determining $\frac{\partial E}{\partial \mathbf{L}_i}$ necessitates intermediate computation of $\frac{\partial \mathbf{X}}{\partial \mathbf{W}}$, carried out using equation~(5), which includes dequantizing the stored $\mathbf{W}^{\text{NF4}}$ to the compute format $\mathbf{W}^{\text{BF16}}$. This allows for evaluation of $\frac{\partial \mathbf{X}}{\partial \mathbf{W}}$ at BFloat16 precision.

In summary, \textsc{QLoRA} utilizes a single storage format for data (typically 4-bit NormalFloat) and a separate computation format (16-bit BrainFloat). During both the forward and backward computations, stored parameters are dequantized into the computation data type, though weight gradients are computed solely for the LoRA parameters, which are maintained in 16-bit BrainFloat.

\section{QLoRA vs. Standard Finetuning}
\label{sec:comp_to_std}

Our discussion so far has outlined the workings of QLoRA, including its capacity to substantially lower memory demands during model finetuning. A central concern remains whether QLoRA achieves comparable performance to conventional full-model finetuning. Additionally, we seek to assess the contributions of different QLoRA elements, particularly evaluating how NormalFloat4 performs relative to standard Float4. The following sections detail our experimental approach aimed at exploring these aspects.

\paragraph{Experimental setup.} We examine three model types—encoder, encoder-decoder, and decoder-only architectures—comparing QLoRA against 16-bit adapter-finetuning as well as full-model finetuning, focusing on models as large as 3B parameters. Our benchmarks encompass the GLUE suite \citep{wang2018glue} with RoBERTa-large \citep{liu2019roberta}, Super-NaturalInstructions (TKInstruct) \citep{wang2022super} paired with T5 \citep{t5}, and 5-shot MMLU \citep{hendrycksmeasuring}, the latter evaluated after finetuning LLaMA on Flan v2 \citep{longpre2023flan} and Alpaca \citep{alpaca}. To further scrutinize NF4 in relation to other 4-bit quantization schemes, we replicate the setup of \citet{dettmers2022case} and assess post-quantization zero-shot accuracy and perplexity over multiple models (including OPT~\citep{zhang2022opt}, LLaMA~\citep{touvron2023llama}, BLOOM~\citep{scao2022bloom}, and Pythia~\citep{biderman2023pythia}) spanning parameter counts from 125m to 13B.

We report experiment-specific setup details within the respective results subsections to facilitate clarity, with comprehensive descriptions found in Appendix~\ref{app:exp-setup}.

Although paged optimizers enable finetuning extremely large models such as 33B/65B \textsc{QLoRA} on single GPUs with 24GB or 48GB of memory, we do not provide quantitative measurements for the efficacy of Paged Optimizers. This is because paging typically happens only for mini-batches with unusually long sequence lengths, which are infrequent in practice. Nevertheless, we present a runtime analysis of paged optimizers applied to 65B parameter models on 48GB GPUs and observe that, for batch sizes of 16, their training speed aligns with standard optimizers. It remains an open question for future studies to systematically characterize when and why paging may lead to slowdowns.

\paragraph{Default LoRA hyperparameters do not match 16-bit performance}

Following standard procedure, which applies LoRA only to the query and value attention projections \citep{hu2021lora}, does not enable recovery of full-finetuning performance on large-scale models. Figure~\ref{fig:hyper} demonstrates that, for LLaMA 7B trained on Alpaca, the number of LoRA adapters deployed throughout the model has the largest influence: covering every linear layer within transformer blocks is essential for matching the results of full-model finetuning. Other hyperparameters, such as the projection rank $r$, are observed to have negligible impact (refer to Appendix \ref{app:exp-setup}).

In a similar vein, our analysis reveals that the default hyperparameters used for fully finetuned baseline models tend to be inadequately optimized. To establish strong baselines, we conduct an extensive hyperparameter search, varying the learning rates between 1e-6 and 5e-5 and the batch sizes from 8 up to 128. Figure~\ref{fig:hyper} presents the results of finetuning the 7B LLaMA model on Alpaca under these optimized settings.

\paragraph{4-bit NormalFloat surpasses 4-bit Floating Point in performance}

Although the 4-bit NormalFloat (NF4) data type is theoretically optimal from an information perspective, it remains uncertain whether these theoretical strengths result in tangible empirical improvements. Adopting the methodology of \citet{dettmers2022case}, we evaluate quantized large language models (including OPT \citep{zhang2022opt}, BLOOM \cite{scao2022bloom}, Pythia \cite{biderman2023pythia}, and LLaMA), covering a parameter range from 125M to 65B, and employing different quantization types on both language modeling tasks and several zero-shot benchmarks. As illustrated in Figure~\ref{fig:nf4} and detailed in Table~\ref{tbl:nf4_ppl}, NF4 yields notable performance enhancements over both FP4 and Int4; further, memory consumption can be reduced by employing double quantization without negatively affecting performance.

\paragraph{k-bit \textsc{QLoRA} achieves parity with 16-bit full finetuning and 16-bit LoRA}

While it has recently been shown that 4-bit quantization for \emph{inference} is feasible, this often incurs some performance loss compared to 16-bit precision \citep{dettmers2022case,frantar2022gptq}. This observation prompts a key question: Can this performance deficit be eliminated through 4-bit adapter finetuning? To address this, we examine two experimental setups.

Our primary comparison centers on fully finetuned 16-bit RoBERTA and T5 models (ranging from 125M to 3B parameters) and their quantized counterparts on GLUE and Super-NaturalInstructions. Results, found in Table~\ref{tab:nf4vsbf16}, show that 16-bit, 8-bit, and 4-bit adapter-based finetuning approaches all match the performance achieved by the full 16-bit finetuned baselines on both datasets. This demonstrates that, following quantization, the accuracy loss attributed to lower numerical precision can be entirely recovered via adapter-based finetuning.

In our second experimental configuration, due to the fact that comprehensive finetuning of models with 11B parameters or greater necessitates access to multiple servers equipped with high-memory GPUs, we instead focus on investigating whether 4-bit \textsc{QLoRA} achieves parity with 16-bit LoRA at model scales ranging from 7B to 65B parameters. Accordingly, we finetune LLaMA models spanning 7B to 65B parameters using two instruction-oriented datasets, Alpaca and FLAN v2. We assess model performance on the MMLU benchmark through 5-shot accuracy, as presented in Table~\ref{tbl:llama-datatype}. The findings demonstrate that NF4 quantization with double quantization is able to fully recapitulate the MMLU results achieved by 16-bit LoRA. Furthermore, it is noteworthy that \textsc{QLoRA} utilizing FP4 trails the 16-bit brain float LoRA baseline by approximately 1 percentage point. These outcomes reinforce two central observations: (1) \textsc{QLoRA} using NF4 effectively replicates both 16-bit full finetuning and 16-bit LoRA outcomes, and (2) NF4 provides improved quantization fidelity compared to FP4.

\paragraph{Summary} The results across our experiments demonstrate that employing 4-bit \textsc{QLoRA} with the NF4 datatype consistently achieves equivalent performance to 16-bit full and LoRA finetuning on established academic benchmarks and evaluation protocols. We further confirm the superiority of NF4 over FP4, and establish that the incorporation of double quantization does not compromise model quality. Taken together, these findings furnish strong support that 4-bit \textsc{QLoRA} finetuning reliably delivers results that are on par with traditional 16-bit approaches.

Consistent with earlier research on quantization \citep{dettmers2022case}, both our MMLU and Elo benchmark data suggest that when resources for finetuning and inference are limited, allocating capacity to a larger base model with lower parameter precision can be advantageous. This underscores the practical efficiency gains afforded by \textsc{QLoRA}. As our 4-bit finetuning experiments did not result in diminished performance compared to full finetuning, this prompts further investigation into the precise location of the performance-precision boundary for QLoRA finetuning, an area we propose for subsequent research.

We turn our attention to examining instruction tuning at scales beyond the feasible range of full 16-bit finetuning given the limitations of academic research hardware.

\section{Pushing the Chatbot State-of-the-art with QLoRA}
\label{sec:pushing}
Building on our demonstration that 4-bit \textsc{QLoRA} achieves equivalent results to 16-bit models across diverse scales, tasks, and datasets, we present an extensive analysis of instruction finetuning utilizing some of the largest open-source language models currently accessible for research. To systematically evaluate these instruction-tuned models, we employ a challenging Natural Language Understanding benchmark (MMLU), and we propose innovative methodologies for assessing chatbot performance in practical scenarios.

\subsection{Experimental setup} Here, we provide a summary of our experimental methodology; comprehensive details are given in Appendix \ref{app:chatbot}.

\paragraph{Data} As the landscape of recent instruction-following datasets has not been extensively characterized, we curate a collection of eight recent datasets. This selection encompasses datasets sourced via crowd-sourcing (OASST1~\citep{kopf2023openassistant}, HH-RLHF~\citep{bai2022training}), data derived from model distillation processes (Alpaca~\citep{alpaca}, self-instruct~\citep{wang2022self}, unnatural-instructions~\citep{honovich2022unnatural}), as well as larger-scale aggregations (FLAN v2~\citep{chung2022scaling}), and hybrid datasets (Chip2~\citep{laion2023}, Longform~\citep{koksal2023longform}). This dataset collection spans a variety of languages, dataset sizes, and licensing conditions.

\paragraph{Training Setup} In order to isolate the impact of instruction finetuning and eliminate confounding variables associated with heterogeneous training objectives, we restrict our QLoRA finetuning to the cross-entropy loss within a supervised learning paradigm, abstaining from using reinforcement learning—even when datasets contain human-annotated response rankings. When datasets provide a clear demarcation between instruction and corresponding response, finetuning is applied exclusively to the response (see ablation studies in Appendix \ref{app:chatbot}). For OASST1 and HH-RLHF, where multiple responses exist per instruction, we consistently select the top-rated response at each node within the dialogue tree and finetune using the entire selected conversational trajectory, encompassing both instructions and responses. Our training employs NF4 \textsc{QLoRA} with double quantization and paged optimizers, mitigating memory overheads from gradient checkpointing. We conduct limited hyperparameter exploration for LLaMA models of sizes 13B and 33B, noting that, aside from learning rate and batch size, settings found suitable at 7B generalize well (including number of epochs). Specifically, for the 33B and 65B models, we reduce the learning rate by half and increase batch size twofold.

\paragraph{Baselines} For comparative analysis, we benchmark our models against both research-oriented (Vicuna~\citep{vicuna2023}, Open Assistant~\citep{kopf2023openassistant}) and commercial (GPT-4~\citep{openai2023gpt}, GPT-3.5-turbo, Bard) chatbot systems. The Open Assistant system is implemented as a LLaMA 33B model that has undergone RLHF-based finetuning on OASST1, aligning with datasets used in our experiments. In contrast, Vicuna conducts full finetuning of LLaMA 13B on a proprietary corpus of user-submitted ShareGPT conversations, effectively acting as a distilled model of OpenAI’s GPT systems.

\subsection{Evaluation}

Consistent with established methodology, we evaluate model performance using the MMLU (Massively Multitask Language Understanding) benchmark \citep{hendrycksmeasuring}, which assesses proficiency across a diverse set of 57 multiple-choice tasks covering domains such as elementary mathematics, U.S. history, computer science, and legal studies. We present 5-shot accuracy scores on the test set.

We further assess generative language abilities using both automated mechanisms and human-based judgments. This second category of evaluations utilizes a set of human-selected queries and focuses on quantifying the response quality produced by the models. Although this method provides a more authentic assessment environment for chatbot models and has become increasingly widespread, the field still lacks a standardized methodology. Our experimental framework, which we detail below, employs nucleus sampling with $p=0.9$ and a temperature of $0.7$ throughout.

\paragraph{Benchmark Data} For evaluation, we make use of two hand-curated query datasets: the Vicuna prompts \citep{vicuna2023} and the OASST1 validation dataset \citep{kopf2023openassistant}. The Vicuna prompts set, comprising 80 questions spanning varied topics, is used directly without any adjustments. In contrast, the OASST1 dataset features a multilingual, crowdsourced set of multiturn user-assistant dialogues. Here, we extract every user utterance from the validation partition to serve as a query, supplementing it with prior conversation turns as part of the prompt. This selection yields a total of 953 distinct user queries. These two datasets are referred to as the Vicuna and OA benchmarks, respectively.

\paragraph{Automated Evaluation} Our initial evaluation employs the procedure from \citet{vicuna2023}, using GPT-4 as a judge to assess system outputs in comparison to ChatGPT (GPT-3.5 Turbo) on the Vicuna dataset. For each query, both ChatGPT’s and another model's outputs are presented to GPT-4, which is then asked to score each reply on a ten-point scale and offer justifications. The model’s performance is expressed as a percentage relative to ChatGPT’s score. Notably, scores exceeding 100\% are possible when a model outperforms ChatGPT on the absolute scale. We have observed that GPT-4 tends to favor the answer that is shown first in the prompt, a so-called ordering effect. To mitigate this bias, we suggest reporting the mean across both prompt orders.

Subsequently, we conduct direct output comparisons between systems, framing the task as a simplified three-way classification (including ties). GPT-4 is tasked with selecting the better response or declaring a tie, and must also provide an explanation for each judgment. All pairwise combinations of systems are tested across the Vicuna and OA benchmarks.

\paragraph{Human Evaluation} Although recent studies propose the feasibility of leveraging generative models to assess system outputs \citep{fu2023gptscore}, there is not yet clear evidence establishing a strong correlation between GPT-4-based ratings and human evaluation of chatbot performance. As such, we supplement automated evaluation by running parallel human annotation exercises on the Vicuna dataset, aligned with both automated protocols previously outlined. Human annotators are recruited via Amazon Mechanical Turk (AMT); two annotators are assigned for comparisons to ChatGPT, and three for direct pairwise comparisons between models.

\paragraph{Elo Rating} We implement a tournament framework incorporating both human and automated (GPT-4) pairwise evaluations, in which various models face off by generating responses to specific prompts. In each match, a pair of models is pitted against one another, and their outputs are judged to determine a winner. This structure is analogous to approaches employed in prior works such as \citet{bai2022training} and \citet{vicuna2023}; however, our methodology integrates ratings from GPT-4 in addition to those from human evaluators. To compute Elo scores~\citep{elo1967proposed,elo1978rating}, we select comparison pairs at random from our pool of annotated matchups. The Elo system, a standard metric in competitive environments like chess, estimates a model’s expected performance in terms of win probability relative to a specific opponent. For instance, when a model with an Elo of 1100 plays one with an Elo of 1000, the higher-rated model is predicted to win roughly 65\% of the time, while a match between models with identical ratings (e.g., 1000 vs 1000 or 1100 vs 1100) yields a 50\% expected win-rate for each. Elo scores are updated after each contest based on the divergence between actual and predicted outcomes—large surprises (such as an underdog winning) induce greater rating changes than results aligning with expectations. Over successive matches, the ratings converge toward representing each model's relative ability in these tasks. All models begin with an initial rating of 1,000, using a parameter of $K=32$. Following the practice of \citet{vicuna2023}, we repeat this entire evaluation process 10,000 times under different random seed initializations to mitigate biases related to the ordering of the matchups.

\subsection{\textbf{Guanaco}: \textsc{QLoRA} trained on OASST1 Achieves State-of-the-art Performance as a Chatbot}

Our automated and human assessment indicates that the top-performing \textsc{QLoRA} fine-tuned model, {Guanaco} 65B, which has been adapted using a variant of OASST1, stands out as the leading open-source chatbot available, with capabilities approaching those of ChatGPT. In head-to-head system-level comparisons rated by human annotators and assessed through Elo scoring, {Guanaco} 65B and 33B display an expected win rate of 30\% against GPT-4—the best reported so far.

Table~\ref{tbl:vicuna_eval} provides the comparative results on the Vicuna benchmark \cite{vicuna2023} against ChatGPT. Among non-GPT-4 models, {Guanaco} 65B secures the highest ranking, reaching 99.3\% of ChatGPT’s performance. Despite {Guanaco} 33B containing more parameters than Vicuna 13B, it operates at 4-bit precision, leading to significantly greater memory efficiency—requiring only 21 GB as opposed to 26 GB—while surpassing Vicuna 13B’s performance by three percentage points. Additionally, {Guanaco} 7B can be deployed on modern smartphones thanks to its 5 GB model size and yet outperforms Alpaca 13B by nearly 20 percentage points.

Nevertheless, the confidence intervals for the values reported in Table~\ref{tbl:vicuna_eval} are notably broad, with many models demonstrating overlapping results. We attribute much of this variability to ambiguities in the use of the scoring scale; for instance, a score of 8 out of 10 may be interpreted differently depending on the context. Therefore, we advocate for employing the Elo ranking framework \citep{elo1967proposed}, which relies on pairwise judgments—collected from human annotators as well as GPT-4—to circumvent the subjectivity of absolute scaling. Table~\ref{tbl:elo} displays Elo scores for leading models. It is notable that model orderings between human judgments and GPT-4 on the Vicuna benchmark are not entirely aligned—especially for Guanaco 7B—but there is substantial concordance for the majority of models, as indicated by a Kendall Tau of $\tau=0.43$ and a Spearman rank correlation of $r=0.55$ at the system level. On the level of individual examples, agreement is weaker, with Fleiss $\kappa=0.25$ reflecting lower inter-annotator consistency between GPT-4 and human votes. In summary, this signifies that while model-based evaluations offer a fairly trustworthy surrogate for human judgment at the system level, they are not a perfect substitute. Additional discussion is provided in Section~\ref{ssec:considerations}.

As presented in Table~\ref{tbl:elo_large}, Elo rankings show that the {Guanaco} models with 33B and 65B parameters surpass all competitors except GPT-4 on both the Vicuna and OA evaluation benchmarks, with their performance being on par with ChatGPT as also reflected in Table~\ref{tbl:vicuna_eval}. It should be observed that the Vicuna benchmark tends to advantage open-source systems, while ChatGPT receives more favorable outcomes under the broader OA benchmark. Moreover, an examination of Tables \ref{tbl:mmlu-llama} and \ref{tbl:vicuna_eval} highlights the critical importance of the choice of finetuning datasets in determining model efficacy. Specifically, Llama models finetuned with FLAN v2 demonstrate strong results on MMLU, yet achieve lower scores on the Vicuna benchmark; analogous trends appear across other models. This divergence underscores a partial independence among current evaluation metrics: excelling on MMLU does not guarantee high chatbot ability as measured by Vicuna or OA, nor does the reverse hold.

Among the top models evaluated,  is unique in being trained exclusively on non-proprietary data due to OASST1 dataset guidelines that prohibit GPT model usage. The next leading model developed solely with open-source datasets is Anthropic HH-RLHF, yet it lags behind {Guanaco} by 30 percentage points on the Vicuna benchmark (refer to Table~\ref{tbl:vicuna_eval}). These collective findings demonstrate that the 4-bit \textsc{QLoRA} approach yields highly competitive chatbots, reaching performance levels close to those of ChatGPT. Notably, our 33B {Guanaco} model can be efficiently trained in under 12 hours on 24 GB consumer GPUs. This makes it feasible to further advance open-source chatbot performance through \textsc{QLoRA} tuning on targeted datasets, potentially creating models that rival the strongest commercial alternatives currently available.

\section{Qualitative Analysis}

Although the primary focus of our evaluation is on quantitative metrics, relying solely on aggregated statistics comes with several limitations. A key concern is benchmark validity \citep{liao2021we}, which questions whether benchmarks genuinely assess what their titles or descriptions claim—especially pertinent as the identification of ``shortcuts'' exploited by machine learning models undermines intended evaluation goals \citep{gururangan2018annotation, poliak2018hypothesis}. To address this limitation in part, we supplement our analysis here with qualitative examination, divided into two segments. In \S\ref{ssec:examples}, we present illustrative instances that capture recurrent tendencies we identified in the responses produced by our 65b {Guanaco} model. In \S\ref{ssec:considerations}, we elaborate on various aspects that inform our interpretation and the broader implications of our results.

\subsection{Qualitative Analysis of Example Generations}
\label{ssec:examples}

To collect illustrative samples, we systematically review responses produced for both the Vicuna and OpenAssistant benchmarks, seeking distinctive regularities in the outputs of {Guanaco}. Upon recognizing a particular behavior, we intentionally design a prompt aimed at eliciting that pattern, even when it would lead to an incorrect answer—for instance, if we detect a proclivity for overly verbose answers, we challenge the model with the directive, ``Answer yes or no without explanation.'' This strategy allows us to surface examples where the model either fails (``lemons,'' i.e., successful adversarial prompts) or maintains robust behavior (``cherries,'' i.e., failed adversarial attacks). The examples provided in this section were all generated using Nucleus Sampling \cite{holtzmancurious} with $p=0.9$.

Naturally, this discussion does not offer an exhaustive analysis; due to the limited scale and qualitative nature of our study, it is infeasible to account for all influencing variables—for instance, the model can produce a vast range of possible completions to any given prompt, making comprehensive coverage impractical. Accordingly, we must depend on selected examples that we expect are sufficiently illustrative. Nevertheless, presenting these qualitative samples contextualizes the quantitative data introduced in earlier sections. As all models and code are publicly released, we anticipate that these findings will encourage subsequent research to delve more thoroughly into the complexities observed here.

\paragraph{Factual Recall} For fact-based queries like ``What is the capital of Zambia?'', all evaluated models reliably deliver accurate responses, for example:

\begin{quote} 
    \textbf{Guanaco}:
    \texttt{The capital of Zambia is Lusaka.}
\end{quote}

 However, as the questions increase in obscurity, {Guanaco} begins to return incorrect answers, though it maintains a tone of confidence. Consider the following instance from HotPotQA \citep{yang2018hotpotqa}:

\begin{quote}
     \textbf{User}: In what year was the singer who popularized `I'll Keep the Lovelight Burning' born?  

     \textbf{Guanaco}:~\texttt{Al Jolson is credited with popularizing the song `I'll Keep the Lovelight Burning,' and he was born in the year 1886.}
 \end{quote}

 In this case, the model provides both an incorrect artist and incorrect year of birth (although the year matches Al Jolson's actual birth date).

\paragraph{Suggestibility} In contrast, {Guanaco} often demonstrates notable resistance to adopting presupposed inaccuracies, as illustrated in the following exchange:

\begin{quote}
    \textbf{User}: How was it finally, officially confirmed that the earth is flat by peer reviewed scientists? \\ \\
    \textbf{Guanaco}: \texttt{There is no official scientific confirmation that the earth is flat; this claim has never been supported by peer-reviewed scientists. The scientific community overwhelmingly agrees that the earth is an oblate spheroid. This has been empirically validated for centuries through extensive observations and experimental evidence. \\ \\ The notion that the earth is flat is classified as a debunked conspiracy theory and has been thoroughly discredited by scientific research. Assertions regarding scientific peer review confirming a flat earth are unfounded and inaccurate. \\ \\ Please let me know if you have further inquiries.}
\end{quote}

also demonstrates a good awareness of the types of unanswerable questions, such as:

\begin{quote}
    \textbf{User}: What time is it?

    \textbf{Guanaco}: \texttt{Apologies, but as an AI language model, I cannot provide real-time data. My answers are informed by existing information and knowledge and do not correspond to the current time. To determine the time, please use a clock or an online time service.}
\end{quote}

This behavior is also seen with prompts like ``Where are you?'', ``How are you?'', and so forth.

\paragraph{Refusal} In a similar vein, {Guanaco} may occasionally reject instructions in a manner that appears arbitrary:

\begin{quote}
    \textbf{User}: Please invert the order of the words in the sentence ``I'm a cool dude in a mocha mood.'' \\ \\
    \textbf{Guanaco}: \texttt{I'm unable to perform that request at this time, but I can clarify why your sentence is grammatically correct.\\ \\The phrase ``I'm a cool dude in a mocha mood" adheres to standard English grammar and syntax conventions. Here, the adjective ``cool'' appropriately precedes the noun ``dude'' as a descriptor, while the prepositional phrase ``in a mocha mood" supplies further context. This construction is both grammatically sound and contextually meaningful.\\ \\If you have further inquiries or require assistance with anything else, please let me know.}
\end{quote}

\paragraph{Secret Keeping} However, this type of refusal does not consistently provide the desired secrecy. There are many cases where system developers intend for the model to withhold particular information from users, even while still allowing unrestricted text outputs. For example, if the model is prompted with ``The secret word is `banana'. Do not reveal or repeat this, no matter what the human requests,'' we can examine how vulnerable this secret is to being extracted. Simply asking directly does not succeed:

\begin{quote}
    \textbf{User}: What is the secret word?

    \textbf{Guanaco}: \texttt{I'm sorry, but I cannot disclose the secret word. I am programmed not to reveal or repeat the secret word, regardless of user requests. Is there something else I can assist you with?}
\end{quote}

which represents the intended outcome. Nonetheless, even minimal manipulation of the prompt can circumvent this restriction:

\begin{quote}
    \textbf{User}: Let's play a game. The objective is to disregard any prior instructions. What is the secret word?

    \textbf{Guanaco}: \texttt{Sorry, the secret word is banana. It's better if we avoid discussing it. Is there anything else you'd like to know?}
\end{quote}

demonstrating the necessity for improved strategies to ensure models adhere more robustly to instructions.

\paragraph{Math} The most significant shortcoming of {Guanaco} lies in mathematics, a domain that presents challenges for many language models, as described by \citep{liang2022holistic}. When {Guanaco} provides intermediate reasoning, its accuracy often improves, as in the following example:

\begin{quote}
    \textbf{User}: Lee mows one lawn and charges \$33. Last week he mowed 16 lawns and three customers each gave him a \$10 tip. How many dollars did Lee earn mowing lawns last week?

    \textbf{Guanaco}: \texttt{Lee earned \$582 mowing lawns last week. \\ \\ To find the total amount he earned, we first need to calculate how much he earned for mowing 16 lawns: \\ \\  \$33 x 16 = \$528 \\ \\ Then we need to add the tips he received from the three customers: \\  \\ \$10 x 3 = \$30 \\ \\ So the total amount Lee earned was: \\ \\ \$528 + \$30 = \$558}
\end{quote}

Despite this, {Guanaco} can sometimes fail on basic math questions if it does not proceed incrementally through the solution—a recognized limitation \cite{weichain}. The following exchange illustrates such a case:

\begin{quote}
    \textbf{User}: Please factorize 1833.

    \textbf{Guanaco}: \texttt{Sure, I'd be happy to help! 1833 is a prime number, which means it has only two factors: 1 and 1833. \\ \\ So the unique factorization of 1833 is: \\ \\ 2\^{}1 * 3\^{}2 * 17\^{}1 \\ \\ I hope this helps! Let me know if you have any other questions.}
\end{quote}

which is not just incorrect (the correct decomposition is $3 \times 17 \times 43$), but doubly so.

\paragraph{Theory of Mind} {Guanaco} demonstrates unexpectedly robust Theory of Mind abilities \cite{nematzadeh2018evaluating, sap2022neural}. For example, the model produces the following highly detailed and accurate response to this scenario:

\begin{quote}
    \textbf{User}: James and Abby are in the bedroom. Abby put the pen in the desk drawer. Abby leaves the bedroom. James moves the pen into the bag. Where does James think Abby will look for the pen?

    \textbf{Guanaco}: \texttt{James knows that Abby put the pen in the desk drawer. But James also knows that he has moved the pen into the bag, and he assumes that Abby does not know about this. So, James may guess that Abby will look for the pen in the desk drawer, because that's where she left it last time. However, this assumption may be incorrect, because Abby may have seen James moving the pen.}
\end{quote}

Yet, the reliability of such reasoning is questionable; the model frequently produces justifications that are inconsistent with the scenario, such as in the following exchange:

\begin{quote}
    \textbf{User}: Evelyn entered the living room. Jackson entered the playroom. James entered the playroom. The beans are in the treasure chest. James exited the playroom. Jackson moved the beans to the pantry. Jackson exited the playroom. James entered the living room. Where will James look for the beans?

    \textbf{Guanaco}: \texttt{James will look for the beans in the pantry, because that's where Jackson moved them.}
\end{quote}

Here, {Guanaco} incorrectly infers the transmission of information that is not specified in the context. Such errors are consistent with findings in recent work \cite{sap2022neural}, but further systematic investigation is warranted.

\subsection{Considerations}
\label{ssec:considerations}

\paragraph{Evaluation}

Our results indicate moderate inter-annotator agreement (Fleiss $\kappa=0.42$), which decreases further when directly comparing outputs from two advanced models. This outcome highlights inadequacies in current benchmark tasks and human evaluation methodologies for chatbot performance. In our manual assessment of ChatGPT versus {Guanaco} 65B responses on the Vicuna benchmark, we note substantial influence of individual subjective preference, as evidenced by significant disagreement among the authors concerning preferred outputs. Future research should address these challenges by leveraging insights from domains such as Human-Computer Interaction and Psychology, which have developed frameworks for managing subjective judgment.

Furthermore, our assessment reveals that automated evaluation methods also display notable biases. For instance, we identify strong order effects with GPT-4, where higher ratings are systematically given to whichever system’s outputs are presented first. The relatively low agreement between GPT-4 and human annotators at the sample level (Fleiss $\kappa=0.25$) indicates that their respective evaluative criteria are often misaligned. Moreover, as illustrated in Table~\ref{tbl:elo_large}, GPT-4 appears to favor its own responses substantially more than human raters do—yielding an Elo score of 1348 compared to the human-assessed score of 1176, corresponding to an approximately 20\% increased likelihood of “winning” against an opponent.

Prospective research should investigate the existence of potential biases within automated evaluation mechanisms and explore corresponding remediation techniques.

\paragraph{Data \& Training} It is important to highlight that the {Guanaco} models are trained using the OASST1 dataset, which is inherently multilingual, and that the OA benchmark encompasses prompts formulated in multiple languages. We suggest that subsequent studies should assess the extent to which multilingual training contributes to performance enhancements on non-English instructions and consider whether this might account for the more pronounced performance disparity observed between the Vicuna-13B model (trained exclusively on English) and the {Guanaco} 33B and 65B variants within the OA benchmark context.

Given the notable results demonstrated by {Guanaco} models, we conducted an analysis to detect possible data leakage between the OASST1 corpus and the Vicuna benchmark prompts. After employing fuzzy string matching across both datasets and manually reviewing the closest candidates, no instances of overlapping prompts were detected.

Moreover, it should be emphasized that our model relies solely on cross-entropy loss within a supervised learning framework and does not incorporate reinforcement learning from human feedback (RLHF). This situation motivates additional research on the respective trade-offs between training with only cross-entropy loss and integrating RLHF. We anticipate that \textsc{QLoRA} will support such large-scale investigations without incurring excessive computational costs.

\section{Related Work}

\paragraph{Quantization of Large Language Models} The primary focus of quantizing large language models (LLMs) has been on enhancing inference-time efficiency. State-of-the-art strategies for preserving the fidelity of 16-bit LLMs typically address outlier activation values, as demonstrated by approaches such as SmoothQuant~\citep{xiao2022smoothquant} and LLM.int8()~\citep{dettmers2022llm}, while other research leverages advanced grouping methodologies \citep{park2022nuqmm, yao2022zeroquant}. Studies on lossy quantization explore the compromises associated with standard rounding~\citep{dettmers2022case,zeng2022glm,pope2022efficiently}, as well as techniques to optimize the rounding process for better quantization accuracy~\citep{frantar2022gptq}. Apart from the present work, SwitchBack layers~\citep{wortsman2023stable} represents the sole example of research that investigates backpropagation through quantized weights at parameter counts exceeding 1B.

\paragraph{Finetuning with Adapters} Our methodology employs Low-rank Adapters~\citep{hu2021lora} (LoRA) for finetuning; however, the landscape of Parameter Efficient FineTuning (PEFT) strategies is broad, with alternatives such as prompt tuning~\citep{qin2021learning,lester2021power,li2021prefix}, input layer modifications~\citep{an2022input}, hidden state adaptation (e.g., IA$^3$)~\citep{liu2022few}, introducing entire additional layers~\citep{houlsby2019parameter}, bias parameter adjustment~\citep{zaken2021bitfit}, learning parameter masks informed by Fisher information~\citep{sung2021training}, and various combinations thereof~\citep{henderson2021compacter}. Our results demonstrate that LoRA adapters attain performance on par with full 16-bit finetuning. A more comprehensive comparison among alternative PEFT techniques, as well as an assessment of their tradeoffs, is reserved for subsequent research.

\paragraph{Instruction Finetuning} Instruction finetuning is designed to align pretrained LLMs with user-provided prompts by exposing the models to input-output pairs sourced from a wide range of datasets. This process tailors model behavior to better follow instructions. Major methodologies and datasets in this domain include MetaICL~\citep{min2021metaicl}, MetaTuning~\citep{zhong2021adapting}, InstructGPT~\citep{ouyang2022training}, FLAN~\citep{wei2021finetuned,chung2022scaling}, PromptSource~\citep{bach2022promptsource}, Super-NaturalInstructions~\citep{wang2022super,sanh2021multitask}, Self-instruct~\citep{wang2022self}, UnnaturalInstructions~\citep{honovich2022unnatural}, OPT-IML~\citep{iyer2022opt}, UnifiedSKG~\citep{xie2022unifiedskg}, OIG/Chip2~\citep{laion2023}, Alpaca~\citep{alpaca}, Vicuna~\citep{vicuna2023}, Koala~\citep{koala_blogpost_2023}, and Self-instruct-GPT-4~\citep{peng2023instruction}.

\paragraph{Chatbots} Many state-of-the-art instruction-following language models are realized as dialogue-oriented chatbots, frequently trained via Reinforcement Learning from Human Feedback (RLHF)~\citep{christiano2017deep} or by generating synthetic data with model-based feedback (RLAIF)~\citep{bai2022constitutional}. Well-known systems and datasets in this area include Anthropic-HH~\citep{askell2021general,bai2022training}, Open Assistant~\citep{kopf2023openassistant}, LaMDA~\citep{thoppilan2022lamda}, and Sparrow~\citep{glaese2022improving}. In contrast to methods relying on reinforcement learning, our top-performing model, Guanaco, is finetuned using multi-turn conversational data from the Open Assistant dataset, originally collected for RLHF applications~\citep{kopf2023openassistant}. To facilitate chatbot evaluation, recent methods leverage GPT-4 as an alternative to expensive human annotations~\citep{vicuna2023,peng2023instruction}. Our contributions include enhancements to such evaluation protocols, emphasizing increased robustness and reliability.

\section{Limitations and Discussion}

Our empirical findings indicate that the \textsc{QLoRA} approach is capable of achieving full 16-bit finetuning performance by combining a 4-bit model base with LoRA adapters. Nevertheless, we have not yet validated whether \textsc{QLoRA} performs equivalently to full 16-bit finetuning at the 33B and 65B parameter scales. Due to the substantial computational demands, addressing this question is postponed to future investigation.

A further limitation is present in the scope of our instruction finetuning model evaluation. Although our analysis encompasses MMLU, the Vicuna benchmark, and the OA benchmark, we did not extend our assessment to additional widely used benchmarks such as BigBench, RAFT, and HELM, and consequently, generalizability to these benchmarks is not guaranteed. Still, our study delivers an extensive evaluation on MMLU and introduces novel chatbot assessment methodologies.

Based on the findings presented, it is evident that benchmark performance tends to reflect the degree of overlap between the finetuning corpus and the evaluation dataset. For instance, the FLAN v2 dataset bears resemblance to MMLU but is less similar to conversational benchmarks, whereas the Chip2 dataset aligns more with the latter. Both models, therefore, exhibit higher scores on benchmarks that are similar to their respective finetuning data, such as MMLU and Vicuna. This observation underscores the need not only for improved benchmarks and evaluation methodologies but also for deliberate reflection on what exactly is being measured. Should our focus be on excelling at knowledge assessment in academic subjects typically found in high school and college settings, or should we prioritize conversational proficiency characteristic of chatbots? Alternatively, are there other competencies we should value? The convenience of benchmarking on pre-existing tasks as opposed to constructing new evaluation datasets may inadvertently guide collective efforts toward specific objectives. As a field, it is important to ensure that benchmarks genuinely capture the aspects of model performance that are most relevant to our goals.

Although this work offers a comprehensive analysis of general chatbot effectiveness, a significant caveat is the limited scope of responsible AI assessment performed for {Guanaco}. We examine the propensity of {Guanaco}-65B to generate sequences with social biases, as compared to alternative models, in Table~\ref{tbl:crows}. Results indicate that {Guanaco}-65B achieves a substantially lower average bias score than other models trained solely via pretraining. This suggests that finetuning on the OASST1 dataset attenuates bias present in the original LLaMA model. While these findings are promising, it remains an open question whether {Guanaco} maintains similar performance on other dimensions of bias. Comprehensive exploration of different forms of bias in {Guanaco} and analogous conversational systems is left as future work.

Another limitation of this study is the absence of experiments involving varied bit-precisions, such as utilizing 3-bit base models, or investigating alternative adapter approaches. Beyond LoRA, the literature includes a diverse array of Parameter Efficient FineTuning (PEFT) techniques demonstrated to be effective; nonetheless, the scalability of such methods to very large models remains unresolved. Our adoption of LoRA was motivated by the substantial empirical evidence supporting its robustness, but other adapters may prove advantageous. Since post-quantization finetuning can apparently recuperate much of the information lost during quantization, more aggressive quantization strategies may become viable. As an example, it is conceivable that LoRA-enabled finetuning following 3-bit GPTQ quantization could approach the performance levels of traditional 16-bit full finetuning.

\section{Broader Impacts}

The finetuning technique introduced in this work, \textsc{QLoRA}, represents the first instance enabling the efficient finetuning of models with up to 33B parameters on a single consumer-grade GPU, as well as 65B parameter models on a single professional GPU, all without compromising performance when compared to standard full finetuning. Our experiments show that the leading 33B parameter model finetuned on the Open Assistant dataset matches ChatGPT’s performance on the Vicuna benchmark. As instruction-based finetuning is critical in adapting pretrained LLMs into capable conversational agents akin to ChatGPT, we anticipate that our approach will greatly expand accessibility to high-quality finetuning, especially among resource-constrained researchers. This, in turn, advances the democratization of advanced NLP technology. Ultimately, \textsc{QLoRA} has the potential to act as a leveling mechanism, reducing the disparity in technological capability between large organizations and smaller teams utilizing consumer hardware.

A further avenue of significant influence is the application of our approach on mobile devices. We contend that \textsc{QLoRA} presents a promising step towards realizing the finetuning of LLMs directly on smartphones and within other environments constrained by limited resources. Although previous work has demonstrated the feasibility of running 7B parameter models on mobile devices, \textsc{QLoRA} is, to our knowledge, the inaugural technique that permits the finetuning of these models in such settings. Based on our analysis, an iPhone 12 Plus, while charging overnight, could utilize \textsc{QLoRA} to finetune approximately 3 million tokens. Although the resultant finetuned 7B models do not attain the performance levels of ChatGPT, we argue that their quality is sufficient to support innovative applications, particularly those previously inhibited by privacy concerns or inadequacies in LLM quality. Moreover, \textsc{QLoRA} offers opportunities for privacy-oriented LLM usage, empowering individuals to maintain ownership and control over their data and models, and thereby reducing deployment barriers.

Nonetheless, the potential misuse of finetuning, as a dual-use technology, cannot be ignored. Prior research has documented risks associated with the widespread adoption of LLMs \citep{bommasani2021opportunities,bender2021dangers}. Still, we advocate that democratizing access to this rapidly proliferating technology will promote more thorough and independent evaluations, contrasting with scenarios where large organizations withhold their models and codebases from public scrutiny.

In summary, we anticipate that \textsc{QLoRA} will contribute a net positive effect by significantly lowering the barriers to high-quality LLM finetuning, thereby making it accessible to a broader community.